% ============================================================================
% RASCUNHO — Seção: Construção do painel analítico definitivo e diagnóstico
% exploratório atualizado. Fragmento LaTeX (requer \usepackage{booktabs}).
% Fonte dos números: construir_painel_definitivo.py e
% diagnostico_painel_definitivo.py, executados sobre a base SIH/SUS 2015-2025.
% ============================================================================

\section*{Construção do painel analítico definitivo}

A aplicação sequencial dos critérios de inclusão e exclusão formalizados pela equipe reduziu a base bruta de \textbf{830 estabelecimentos e 7.016 observações hospital-ano} para um painel balanceado de \textbf{317 hospitais acompanhados ao longo dos onze anos do período (2015--2025), totalizando 3.487 observações hospital-ano}. A primeira etapa, relativa ao tipo de estabelecimento, excluiu 76 unidades --- hospitais dia (código 62), hospitais psiquiátricos (código 07 combinado à especialização 006) e Centros de Atenção Psicossocial (código 70) ---, preservando explicitamente as maternidades (especialização 005). Dado que o tipo declarado varia entre anos para um mesmo CNES, a classificação adotou o \textbf{valor modal} de cada estabelecimento nos anos com produção, com os casos de cadastro instável registrados em log próprio de revisão. A segunda etapa aplicou o critério de porte em caráter fixo: calculada a \textbf{mediana do número total de leitos} de cada hospital ao longo dos anos com produção registrada, foram excluídos integralmente os 349 estabelecimentos com mediana igual ou inferior a 50 leitos --- o filtro de maior impacto quantitativo, respondendo por 42\% da base original. A terceira etapa removeu seis \textbf{hospitais de campanha} criados exclusivamente para o enfrentamento da COVID-19 e ainda remanescentes após os filtros anteriores; a identificação combinou a nomenclatura do estabelecimento e os rótulos assistenciais associados à pandemia, condicionada à ausência de produção fora do biênio 2020--2021, de modo a não excluir indevidamente hospitais regulares temporariamente rebatizados. A quarta etapa expurgou dos indicadores os \textbf{procedimentos específicos da pandemia (código 999)}, que responderam por 6,7\% da produção do biênio 2020--2021, sem remoção de qualquer observação --- a série temporal preserva integralmente os anos pandêmicos. Por fim, a restrição de \textbf{painel balanceado} eliminou 82 estabelecimentos sem produção em todos os onze anos; a auditoria caso a caso indicou que as ausências refletem majoritariamente fenômenos reais --- hospitais regionais estaduais inaugurados no curso do período e desligamentos progressivos do faturamento SIH, como o do Hospital do Servidor Público Municipal a partir de 2022 --- e não falhas de cobertura da base, cuja contagem anual de estabelecimentos permaneceu estável (615 a 687 CNES por ano). A Tabela~\ref{tab:funil} sintetiza o funil de filtros; o efeito de cada etapa, com a relação nominal dos estabelecimentos removidos e o motivo individual de exclusão, encontra-se nos arquivos de auditoria que acompanham o painel.

\begin{table}[htbp]
\centering
\caption{Funil de construção do painel analítico definitivo (SIH/SUS--SP, 2015--2025)}
\label{tab:funil}
\begin{tabular}{lrrrr}
\toprule
Etapa & \multicolumn{2}{c}{Removidos} & \multicolumn{2}{c}{Remanescentes} \\
\cmidrule(lr){2-3}\cmidrule(lr){4-5}
 & CNES & Hosp.-ano & CNES & Hosp.-ano \\
\midrule
Painel bruto & --- & --- & 830 & 7.016 \\
A. Tipo de estabelecimento & 76 & 646 & 754 & 6.370 \\
B. Porte (mediana de leitos $>50$) & 349 & 2.538 & 405 & 3.832 \\
C1. Hospitais de campanha COVID & 6 & 9 & 399 & 3.823 \\
C2. Procedimentos código 999 & 0 & 0 & 399 & 3.823 \\
D. Painel balanceado (11 anos) & 82 & 336 & \textbf{317} & \textbf{3.487} \\
\bottomrule
\end{tabular}

\smallskip
{\footnotesize Nota: a etapa C2 não remove observações; os indicadores oficiais passam a ser calculados sem os procedimentos específicos da pandemia (6,7\% da produção de 2020--2021).}
\end{table}

\section*{Diagnóstico exploratório do painel definitivo}

No painel filtrado, a \textbf{taxa de mortalidade hospitalar mediana é de 5,35\%} (intervalo interquartílico de 2,99\% a 6,97\%), o \textbf{tempo médio de permanência mediano é de 4,42 dias} (3,51 a 5,78) e o \textbf{custo mediano por saída é de R\$~1.064,68} (R\$~732,35 a R\$~1.719,65). A proporção de alta complexidade é fortemente assimétrica, com mediana de apenas 0,17\% e média de 7,2\%, refletindo a concentração desse tipo de produção em poucos estabelecimentos. Cabe registrar que, no recorte de porte adotado, as duas versões da mortalidade --- com e sem a exclusão dos óbitos obstétricos e perinatais --- tornam-se praticamente indistinguíveis (medianas de 5,35\% e 5,34\%), o que simplifica a leitura dos resultados sem prejuízo da manutenção de ambas as séries. As taxas de ocupação medianas situam-se em 73,5\% na internação e 67,0\% na UTI, calculadas segundo o denominador definido e confirmado pela equipe responsável pela base (permanência acumulada no ano dividida por leitos~$\times$~365); \textbf{valores superiores a 100\% são artefato esperado dessa fórmula} --- decorrente, por exemplo, de variações do parque de leitos declarado ao longo do ano --- e não indicam inconsistência da base, devendo apenas ser considerados na especificação dos modelos de fronteira (Beta/ZOIB) previstos para essas proporções.

A comparação entre o painel bruto e o definitivo (\emph{cf.} figura figD05) evidencia que o \textbf{efeito da seleção de hospitais domina amplamente o efeito da remoção dos procedimentos pandêmicos}: a mediana anual típica de mortalidade desloca-se de 4,42\% para 5,31\%, o tempo de permanência de 3,95 para 4,45 dias e o custo por saída de R\$~789,78 para R\$~1.214,60 quando se passa da base bruta ao conjunto dos 317 hospitais, coerentemente com um painel composto por unidades maiores e de maior complexidade; a subsequente exclusão do código 999 é neutra sobre mortalidade e permanência medianas e concentra seu efeito no custo do biênio pandêmico, cuja mediana-das-medianas anuais recua para R\$~1.059,46.

\subsection*{Distribuição por modelo de gestão (\emph{proxy} provisório)}

A distribuição do painel pela classificação assistencial do SIH --- reiterando que se trata de \textbf{\emph{proxy} provisório do modelo de gestão}, que não captura as categorias Autarquia e PPP e mistura esfera administrativa com modelo de gestão na categoria ``Público Municipal'', na pendência do \emph{crosswalk} institucional definitivo --- compreende cinco categorias (Tabela~\ref{tab:gestao}). O grupo \textbf{Filantrópico} concentra 187 dos 317 estabelecimentos (2.057 observações) e apresenta a segunda maior mortalidade mediana (5,45\%) com o menor tempo de permanência (4,11 dias). O grupo \textbf{Público Municipal}, segundo maior contingente com 60 estabelecimentos (660 observações), exibe simultaneamente a \textbf{maior mortalidade mediana (5,64\%) e o menor custo mediano por saída (R\$~961,89)}, com produção de alta complexidade praticamente nula (0,09\%) --- perfil compatível com hospitais gerais de retaguarda municipal; recorde-se que a permanência dessa categoria no escopo do estudo, bem como a distinção interna entre administração direta municipal e OSS municipal, é questão registrada como pendente nos critérios do projeto. O grupo \textbf{OSS} (39 estabelecimentos, 387 observações) apresenta mortalidade mediana de 4,93\%, permanência de 4,85 dias, custo de R\$~1.147,40 e a maior proporção de alta complexidade entre os grupos públicos (2,42\%), enquanto a \textbf{administração Direta} (33 estabelecimentos, 350 observações) registra a menor mortalidade entre eles (4,60\%) ao lado do maior tempo de permanência do painel (6,18 dias) e custo mediano equivalente ao das OSS (R\$~1.143,88). Os únicos estabelecimentos que transitam entre categorias ao longo do período são os \textbf{cinco casos documentados de conversão Direta$\to$OSS}, sem qualquer outra transição espúria de rótulo --- indício favorável à qualidade do campo no recorte final.

A categoria \textbf{Privado} requer tratamento à parte. O grupo reúne \textbf{três estabelecimentos} (33 observações), de perfil marcadamente atípico --- custo mediano por saída de R\$~12.798,36, proporção de alta complexidade de 98,4\% e mortalidade de 2,56\% ---, e sua composição exige registro explícito: dois prestadores privados contratualizados de nicho (Leforte Liberdade e Unimed Sorocaba) e o \textbf{Hospital Universitário da UFSCar} (CNES 5586348), que, conforme confirmação por fonte externa, é \textbf{hospital público federal administrado pela Ebserh} (Empresa Brasileira de Serviços Hospitalares, vinculada ao Ministério da Educação) --- o rótulo ``Privado'' que o SIH lhe atribui em todos os anos da série \textbf{não corresponde à sua natureza jurídica real}. A equipe decidiu mantê-lo agrupado nessa categoria \textbf{por conveniência estatística}, elevando o grupo de $n=2$ para $n=3$; trata-se de escolha de agrupamento pragmático registrada, e não de endosso da classificação da fonte. A ressalva estatística permanece integralmente: \textbf{com $n=3$, qualquer coeficiente estimado para essa categoria não é interpretável como efeito médio do modelo de gestão}, tanto pelo tamanho amostral quanto pela heterogeneidade interna do grupo, que mistura prestadores privados de nicho com um hospital universitário federal de perfil assistencial distinto.

\begin{table}[htbp]
\centering
\caption{Painel definitivo por modelo de gestão (\emph{proxy} provisório) --- medianas 2015--2025}
\label{tab:gestao}
\begin{tabular}{lrrrrrr}
\toprule
Categoria & CNES & Hosp.-ano & Mort. & TMP (dias) & Custo/saída (R\$) & Alta compl. \\
\midrule
Filantrópico      & 187 & 2.057 & 5,45\% & 4,11 & 1.086,13 & 0,11\% \\
Público Municipal &  60 &   660 & 5,64\% & 4,77 &   961,89 & 0,09\% \\
OSS               &  39 &   387 & 4,93\% & 4,85 & 1.147,40 & 2,42\% \\
Direta            &  33 &   350 & 4,60\% & 6,18 & 1.143,88 & 0,83\% \\
Privado           &   3 &    33 & 2,56\% & 4,24 & 12.798,36 & 98,36\% \\
\bottomrule
\end{tabular}

\smallskip
{\footnotesize Nota: classificação assistencial do SIH utilizada como \emph{proxy} provisório do modelo de gestão --- não captura Autarquia/PPP e mistura esfera administrativa com modelo de gestão; \emph{crosswalk} definitivo pendente. Os cinco hospitais com conversão Direta$\to$OSS figuram em ambas as categorias. A categoria Privado ($n=3$) inclui o Hospital Universitário da UFSCar, hospital público federal administrado pela Ebserh, agrupado aqui por conveniência estatística e não por corresponder à sua natureza jurídica real; o agrupamento eleva o grupo de $n=2$ para $n=3$, mas não resolve a limitação de tamanho amostral discutida no texto, e qualquer leitura do coeficiente dessa categoria deve considerar que ela mistura prestadores privados contratualizados de nicho com um hospital universitário federal de perfil assistencial distinto.}
\end{table}

\subsection*{Complexidade hospitalar}

O painel definitivo distribui-se pelas faixas de complexidade do modelo de Barcelona em 49, 134, 106, 15 e 10 estabelecimentos nas faixas 2 a 6, respectivamente, com três CNES ainda sem classificação, em verificação junto ao responsável pela base. Em cumprimento à decisão da equipe, foram implementadas \textbf{duas versões coexistentes do escore de complexidade}: o \textbf{escore estrutural puro} (pontuação de Barcelona) e o \textbf{escore ponderado por mortalidade} (pontuação multiplicada pela mortalidade relativa do hospital-ano frente à mediana estadual do ano, em fórmula provisória sujeita a ratificação). As duas versões apresentam correlação de Spearman de 0,926 --- fortemente associadas, porém não redundantes. Permanece integralmente vigente a \textbf{salvaguarda de circularidade}: nos modelos em que a mortalidade figura como variável dependente, deve ser utilizado exclusivamente o escore estrutural, reservando-se a versão ponderada aos modelos de permanência, custo, ocupação e produção.

\subsection*{Limitações}

Além da natureza provisória do \emph{proxy} de modelo de gestão, já reiterada, cumpre destacar uma \textbf{limitação de desenho relevante para a estratégia de identificação}: dos cinco hospitais com conversão Direta$\to$OSS documentada no período, \textbf{três (CNES 2082225, 2091755 e 2750511) somente aparecem sob gestão OSS em 2025, oferecendo um único ano de observação pós-tratamento}; os demais convertem-se em 2019 (CNES 2081695, sete anos pós) e 2023 (CNES 2078287, três anos pós). A identificação \emph{within-hospital} do efeito da conversão repousa, portanto, de forma desproporcional sobre dois estabelecimentos, e as estimativas associadas aos conversores tardios devem ser lidas como provisórias, passíveis de fortalecimento apenas com a extensão futura da janela temporal ou com a incorporação, via \emph{crosswalk} institucional, de conversões adicionais hoje não identificáveis --- em particular as envolvendo as categorias Autarquia e PPP, ausentes das fontes disponíveis.

% ============================================================================
% Fim do rascunho.
% DECISÕES REGISTRADAS (não são pendências):
%   - HU-UFSCar: classificação real CONFIRMADA (público federal/Ebserh);
%     mantido agrupado em "Privado" por decisão de agrupamento pragmático
%     da equipe (n=2 -> n=3) — não é erro, é escolha registrada.
%   - Janela de hospitais de campanha: 2020-2021 (verificado que nenhum
%     excluído dependia de 2019).
% Pendências deliberativas sinalizadas no texto:
%   (1) escopo/decomposição de Público Municipal — critérios §6.2;
%   (2) 3 CNES sem faixa Barcelona — Alberto;
%   (3) forma funcional definitiva do escore ponderado por mortalidade — §4;
%   (4) crosswalk institucional definitivo (incl. PPP/Autarquia) — Alberto.
% ============================================================================
